
\documentclass{article}
\usepackage{colortbl}
\usepackage{makecell}
\usepackage{multirow}
\usepackage{supertabular}

\begin{document}

\newcounter{utterance}

\centering \large Interaction Transcript for game `wordle\_withclue', experiment `medium\_frequency\_words\_clue\_no\_critic', episode 5 with selfrefine{-}Qwen3.5{-}9B.
\vspace{24pt}

{ \footnotesize  \setcounter{utterance}{1}
\setlength{\tabcolsep}{0pt}
\begin{supertabular}{c@{$\;$}|p{.15\linewidth}@{}p{.15\linewidth}p{.15\linewidth}p{.15\linewidth}p{.15\linewidth}p{.15\linewidth}}
    \# & \multicolumn{2}{c}{Player} && \multicolumn{2}{c}{Game Master} \\
    \hline

    \theutterance \stepcounter{utterance}  
    & & & \multicolumn{4}{p{0.6\linewidth}}{
        \cellcolor[rgb]{0.9,0.9,0.9}{
            \makecell[{{p{\linewidth}}}]{
                \texttt{\tiny{[P1$\langle$GM]}}
                \texttt{You are a language wizard who likes to guess words by using the given rules.} \\
\\ 
\texttt{Welcome to Wordle! You have six attempts to guess the target word, a valid English word of five lowercase letters (a{-}z). Please use the tags "explanation:" and "guess:" to provide a concise explanation for each guess.} \\
\\ 
\texttt{To help you make an informed guess, you will receive a clue for the word, such as} \\
\texttt{clue: snowy white.} \\
\\ 
\texttt{Here is an example guess based on the clue:} \\
\texttt{explanation: In the fairy tail Snow White, the girl is killed because she eats a poisoned apple. And the word apple has 5 letters.} \\
\texttt{guess: apple} \\
\\ 
\texttt{After each guess, your answer will be validated, and you will receive feedback indicating which letters are correct (green), which letters are correct but in the wrong position (yellow), and which letters are incorrect (red). This feedback can be useful in determining which letters to include or exclude in your next guess.} \\
\\ 
\texttt{For example, the feedback for "apple" might be:} \\
\texttt{guess\_feedback: a<yellow> p<yellow> p<green> l<yellow> e<red>} \\
\\ 
\texttt{The explanation should contain details about how both the guess\_feedback and clue work together in tandem to formulate a new guess.} \\
\\ 
\texttt{Do not generate any other text.} \\
\\ 
\texttt{Let's begin with your first guess.} \\
\\ 
\texttt{clue: tick off} \\
            }
        }
    }
     \\ \\

    \theutterance \stepcounter{utterance}  
    & \multicolumn{4}{p{0.6\linewidth}}{
        \cellcolor[rgb]{0.561,0.737,0.561}{
            \makecell[{{p{\linewidth}}}]{
                \texttt{\tiny{[P1$\rangle$GM]}}
                \texttt{explanation: The clue "tick off" suggests a word related to checking items or marking them, possibly "check". The word "check" has five letters and fits the theme of the clue.} \\
\texttt{guess: check} \\
            }
        }
    }
    & & \\ \\

    \theutterance \stepcounter{utterance}  
    & & & \multicolumn{2}{p{0.3\linewidth}}{
        \cellcolor[rgb]{0.9,0.9,0.9}{
            \makecell[{{p{\linewidth}}}]{
                \texttt{\tiny{[GM$|$GM]}}
                \texttt{attempts = 1} \\
\texttt{target\_word = steam} \\
\texttt{target\_word\_clue = tick off} \\
\texttt{guess = check} \\
\texttt{guess\_feedback = c<red> h<red> e<green> c<red> k<red>} \\
            }
        }
    }
    & & \\ \\

    \theutterance \stepcounter{utterance}  
    & & & \multicolumn{4}{p{0.6\linewidth}}{
        \cellcolor[rgb]{0.9,0.9,0.9}{
            \makecell[{{p{\linewidth}}}]{
                \texttt{\tiny{[P1$\langle$GM]}}
                \texttt{guess\_feedback: c<red> h<red> e<green> c<red> k<red>} \\
\\ 
\texttt{Provide your response only in this format.} \\
\texttt{explanation: details} \\
\texttt{guess: word} \\
            }
        }
    }
     \\ \\

    \theutterance \stepcounter{utterance}  
    & \multicolumn{4}{p{0.6\linewidth}}{
        \cellcolor[rgb]{0.561,0.737,0.561}{
            \makecell[{{p{\linewidth}}}]{
                \texttt{\tiny{[P1$\rangle$GM]}}
                \texttt{explanation: The feedback for "check" confirms only 'e' is correct and in the third position, while 'c', 'h', 'k' are incorrect. The clue "tick off" suggests meanings like "annoy" or "deduct points". Since 'e' is fixed at index 3, words like "score" (s{-}c{-}o{-}r{-}e) are invalid due to 'c'. Words like "stare" (s{-}t{-}a{-}r{-}e) ignore the specific "annoy" meaning. A better approach is to test the root of "annoy" or "deduct". "Anger" fits the "annoy" meaning, has 5 letters, and places 'e' at the end, not the middle, so it conflicts with the 'e' at index 3. Wait, "anger" ends in 'r'. "Annoy" has 5 letters: a{-}n{-}n{-}o{-}y. No 'e'. "Irate"? i{-}r{-}a{-}t{-}e. 'e' is at the end. "Frustrated" is too long. How about "scold"? s{-}c{-}o{-}l{-}d. 'c' is red. "Reprimand" too long. "Scorn"? s{-}c{-}o{-}r{-}n. 'c' is red. "Spite"? s{-}p{-}i{-}t{-}e. 'e' at end. "Gripe"? g{-}r{-}i{-}p{-}e. 'e' at end. "Rant"? 4 letters. "Raves"? r{-}a{-}v{-}e{-}s. 'e' at 4. "Riled"? r{-}i{-}l{-}e{-}d. 'e' at 4. "Fumed"? f{-}u{-}m{-}e{-}d. 'e' at 4. "Sulky"? s{-}u{-}l{-}k{-}y. No 'e'. "Grumpy"? 6 letters. "Badly"? b{-}a{-}d{-}l{-}y. No 'e'. "Mean"? 4 letters. "Nasty"? n{-}a{-}s{-}t{-}y. No 'e'. "Rude"? 4 letters. "Vile"? 4 letters. "Wicked"? w{-}i{-}c{-}k{-}e{-}d. 6 letters. "Evil"? 4 letters. "Sour"? s{-}o{-}u{-}r. 4 letters. "Bitter"? b{-}i{-}t{-}t{-}e{-}r. 6 letters. "Acid"? 4 letters. "Sharp"? s{-}h{-}a{-}r{-}p. 'h' is red. "Sting"? s{-}t{-}i{-}n{-}g. No 'e'. "Prick"? p{-}r{-}i{-}c{-}k. 'c', 'k' red. "Needle"? 6 letters. "Pinch"? p{-}i{-}n{-}c{-}h. 'c', 'h' red. "Nag"? 3 letters. "Nags"? 4 letters. "Naggy"? n{-}a{-}g{-}g{-}y. No 'e'. "Tease"? t} \\
            }
        }
    }
    & & \\ \\

    \theutterance \stepcounter{utterance}  
    & & & \multicolumn{2}{p{0.3\linewidth}}{
        \cellcolor[rgb]{0.9,0.9,0.9}{
            \makecell[{{p{\linewidth}}}]{
                \texttt{\tiny{[GM$|$GM]}}
                \texttt{Error: The response should contain the 'guess:' keyword exactly once.} \\
            }
        }
    }
    & & \\ \\

    \theutterance \stepcounter{utterance}  
    & & & \multicolumn{2}{p{0.3\linewidth}}{
        \cellcolor[rgb]{0.9,0.9,0.9}{
            \makecell[{{p{\linewidth}}}]{
                \texttt{\tiny{[GM$|$GM]}}
                \texttt{game\_result = ABORT} \\
            }
        }
    }
    & & \\ \\

\end{supertabular}
}

\end{document}
